\section{Cusp completeness and the absolute volume}
\label{cg:complete}

Take two oppositely oriented copies of $P_E$. Pair their white
faces by the projection labels. For a pure leaf, pair its shaded
triangles within each copy, fixing $O$ and sending
$C_1\mapsto C_0$, $C_2\mapsto C_3$. For an odd leaf use
\eqref{cg:odd-pairing}. Section~\ref{cg:topology} identifies the
quotient with $Y^\eps_E$, and Section~\ref{cg:maps} proves
smoothness away from the ideal vertices.

At each ideal vertex a horosphere cuts out a Euclidean rectangle.
Its two opposite shaded sides have equal length. A unique choice
of horosphere makes both lengths one. Make this choice independently
at the finitely many ideal vertices. Each shaded ideal triangle
then has three horocyclic corner segments of length one. An
isometry of ideal triangles carrying corresponding vertices to
corresponding vertices preserves precisely this decoration,
including every permutation in \eqref{cg:odd-pairing}.

The white maps preserve the decoration as well: they are the
label-preserving maps between the two identical copies of the
polyhedron, and the shaded-side normalization agrees at each
matched vertex. Thus each cusp link is assembled from Euclidean
rectangles without dilation. The original topological identification
makes it a compact torus. It is a smooth Euclidean torus, since
the edge returns have already excluded cone singularities.

A common shrinking factor makes all chosen horospheres disjoint
and is preserved by all pairings. Beyond each link lies its
complete horoball end. Their complement is compact, so the whole
hyperbolic metric is complete. Equivalently the universal developing
map is a complete local isometry to hyperbolic space and hence
covers it. Every cell and cusp tail has finite volume. All
attachments were unions with disjoint interiors, and the original
cut has degree one. Therefore
\[
 \Vol(Y^\eps_E)=2\Vol(P_E)=2\pi F_0.
 \]
This proves the metric and volume assertions of
Theorem~\ref{cg:main}. The equality is absolute: it follows by
partitioning the actual metric into the displayed cells and using
\eqref{cg:tetra}. No equality of volume derivatives or choice of
an additive regulator constant is involved.

\section{The original crossing-circle curves}
\label{cg:peripheral}

The crossing-circle torus has two rectangles, one from each cut
ball. The disk longitude follows the twice-punctured crossing disk.
The meridian goes once around the auxiliary circle. Both are
marked before the cuts. We trace their words and check
primitiveness against the full Euclidean cusp area.

In the normalized chart \eqref{cg:normalized} these words are
\begin{equation}\label{cg:peripheral-words}
 \mu_0=S^{-1},\qquad \lambda=W_2^{-1}W_0 .
\end{equation}
For $\lambda$, leave one disk side through white face $0$ and
return through face $2$ with its inverse. For $\mu_0$, cross
the shaded identification once in the reverse direction.
The two half-twist collars instead give
\begin{equation}\label{cg:odd-meridians}
 \mu_-=A^{-1}W_0,\qquad \mu_+=W_0^{-1}B .
\end{equation}
The first path passes through white face $0$ and returns through
$A^{-1}$; the second passes through $B$ and returns through
$W_0^{-1}$. These are the two half-disk rotations in the original
disk frame. We label the first negative and the second positive.
Equivalently the retained crossing $R^{-1}$ uses the first collar
and $R$ uses the second. They give the same unmarked triangular
pairing but differ by a full longitude in their meridian paths.
This fixes the sign convention in \eqref{cg:cusp-shape} by
peripheral paths.

Return to \eqref{cg:leaf-coordinates}, where the auxiliary apex
is $\infty$. The normalization is
$M(z)=k(z-1)/(z-E)$ with $k=(1+E)/2$, and $M(\infty)=Z$.
Conjugating the pure words by $M$ gives
\begin{equation}\label{cg:translations}
 \mu_0=E-1,\qquad \lambda=-2(E+1),\qquad
 \mu_-=-2,\qquad \mu_+=2E
\end{equation}
as translations in the original coordinate. For a cross-copy
word the target normalization is the reflected one; composing
the two maps in \eqref{cg:odd-meridians} returns to the original
chart before this conjugation. The matrices
\eqref{cg:white} and \eqref{cg:shaded} give
\eqref{cg:translations} by direct multiplication.
Consequently
\begin{equation}\label{cg:half-shear}
 \mu_\eps=\mu_0-\eps\lambda/2
\end{equation}
as Euclidean vectors on the respective cusp developments.
This compares the markings of different half-twist partners;
it does not make half a longitude an integral curve on the pure
torus.

\begin{figure}[ht]
\centering\begin{tikzpicture}[x=1.25cm,y=1.25cm,>=Latex,every node/.style={font=\small}]
\coordinate(O) at (0,0);
\coordinate(L) at (-4,0);
\coordinate(M) at (0,2);
\fill[blue!8] (-4,0) rectangle (-2,2);
\fill[blue!15] (-2,0) rectangle (0,2);
\draw[thick] (-4,0) rectangle (0,2);
\draw[thick] (-2,0)--(-2,2);
\draw[->,very thick] (O)--(L) node[midway,below=5pt] {$\lambda=-4c$};
\draw[->,very thick] (O)--(M) node[right=5pt] {$\mu_0=2s\,\ii$};
\draw[->,thick,dashed] (O)--(-2,2) node[above=4pt] {$\mu_-$};
\draw[->,thick,dashed] (O)--(2,2) node[above=4pt] {$\mu_+$};
\draw[gray,dotted] (0,2)--(2,2)--(2,0)--(0,0);
\node at (-3,1) {first rectangle};
\node[align=center,text width=1.5cm] at (-1.3,.4) {second\\rectangle};
\node at (-3,-.7) {$2c$};
\node at (-1,-.7) {$2c$};
\fill(O) circle(1.5pt);
\node[below right] at(O) {$0$};
\end{tikzpicture}

\caption{The crossing-circle cusp in the original leaf coordinate,
after rotation by $e^{-\pi\ii a}$. Put $s=\sin\pi a$ and
$c=\cos\pi a$. The pure meridian is $2s\,\ii$ and the disk
longitude is $-4c$. Two rectangles of width $2c$ and height $2s$
give one fundamental parallelogram. The dashed vectors are
$\mu_-=\mu_0+\lambda/2$ and $\mu_+=\mu_0-\lambda/2$,
each for its own half-twist partner.}
\label{cg:fig-cusp}
\end{figure}

The adjacent sides of the original rectangle have lengths
$|E-1|=2\sin\pi a$ and $|E+1|=2\cos\pi a$.
The two rectangles of the cusp have total area
$8\sin\pi a\cos\pi a=4\sin2\pi a$. On the other hand
\begin{equation}\label{cg:lattice-area}
 \left|\operatorname{Im}(\overline{\mu_0}\lambda)\right|
       =4\sin2\pi a .
\end{equation}
The two traced closed curves span a sublattice of the actual cusp
lattice; equality of areas makes its index one. Substitution of
either odd meridian leaves the area unchanged, proving the
primitive marking for those collars as well.

Finally $(E+1)/(E-1)=-\ii\cot\pi a$.
Equations \eqref{cg:translations}--\eqref{cg:half-shear} give
\eqref{cg:cusp-shape}, with
$\operatorname{Im}\tau_\eps=2c_\ell/(1+\eps^2c_\ell^2)>0$.
The bound on $c_\ell$ is the imported leaf window
\eqref{cg:leaf-window}. This completes Theorem~\ref{cg:main}.
All main cusps were already included in the horocycle construction;
no additional eigenvalue or deformation hypothesis is required.

\section{Choosing a knot among the main components}
\label{cg:knot}

The construction kept every main component. A finite endpoint
argument produces a one-component choice throughout the stated
grammar.

\begin{proposition}\label{cg:knot-choice}
One can choose pure or odd leaves in every closed expression of
Definition~\ref{cg:grammar} so that its main strands form one
component. Either sign may be used at an odd leaf. Subsequent
full twists preserve that component count.
\end{proposition}
\begin{proof}
Write the three pairings of four cyclically ordered endpoints as
\[
 I=(0,3)(1,2),\qquad U=(0,1)(2,3),\qquad X=(0,2)(1,3).
\]
A pure leaf has pairing $U$, an odd leaf has pairing $X$, and
turn exchanges $I,U$ and fixes $X$. Tracing the two stack
connections gives
\begin{center}
\begin{tabular}{c|ccc}
lower/upper & $I$ & $U$ & $X$\\ \hline
$I$ & $I$ & $U$ & $X$\\
$U$ & $U$ & $U$ and one internal circle & $U$\\
$X$ & $X$ & $U$ & $I$
\end{tabular}
\end{center}
Only $U\cdot U$ creates an internal component.
The base $hh$ has loop-free possibilities $\{I,U\}$.
Put $A=\{I,U\}$, $B=\{I,X\}$, $C=\{U,X\}$,
$D=\{I,U,X\}$. All products, discarding $U\cdot U$, give
\begin{center}
\begin{tabular}{c|ccc}
 & $A$ & $B$ & $C$\\ \hline
$A$ & $A$ & $D$ & $C$\\
$B$ & $D$ & $B$ & $D$\\
$C$ & $C$ & $D$ & $A$
\end{tabular}
\end{center}
Each result contains at least two choices, also if an input
already contains all three. Attaching a leaf is the row or column
$C$. A turn permutes the possibilities. Inductively every completed
box, and the root stack, has at least two loop-free choices.
At least one is $U$ or $X$, whose cylindrical closure is connected;
the closure of $I$ has two components. Backtracking the choices
specifies the required leaves. Both signs have pairing $X$,
and full twists have even endpoint permutation.
\end{proof}

The last assertion concerns connectivity. Theorem~\ref{cg:main}
proves hyperbolicity of the augmented exteriors, retaining the
crossing circles. Hyperbolicity after filling those circles
requires separate deformation and slope arguments.

\section{Dependencies and finite algebra checks}\label{cg:checks}

The variational existence, uniqueness and leaf interval are the
numbered inputs from \cite{Newton} listed in Section~\ref{cg:intro}.
The cell and complement construction, entire face maps, cusp
decoration and primitive marking are proved here.
The ideal-tetrahedron volume formula \eqref{cg:tetra} is classical
background. The references on augmented links and variational
geometry provide context and do not replace the ordered construction.
No quantum asymptotic enters Theorem~\ref{cg:main}.
Corollary~\ref{cg:source-volume} additionally uses the positive
source limit of \cite[Corollary~5.4]{Newton}.

The deterministic program \texttt{check.py} checks rational-function
identities: nine prism maps, turn and closure compatibility, pure
and odd edge returns, and the complete translations
\eqref{cg:translations}. It enumerates the endpoint pairings
independently of the tables in Section~\ref{cg:knot}. Its negative
controls alter a vertex order and an odd meridian. The output
records source hashes. These finite checks support the displayed
algebra; they do not prove convexity, identify a manifold, establish
completeness or supply the stationary-point theorem.

\begin{thebibliography}{9}
\bibitem{Newton}
\emph{Newton coefficients and positive sources for Conway tangles},
companion paper, Definition~2.1, Proposition~5.3,
Corollary~5.4 and Lemma~6.1.
\href{https://sashakolpakov.github.io/murakamis_dream/papers/newton-tangles.pdf}
{Companion PDF}; source: \texttt{papers/newton-tangles/main.tex}.

\bibitem{Milnor}
J.~Milnor, Hyperbolic geometry: the first 150 years,
\emph{Bull.\ Amer.\ Math.\ Soc.\ (N.S.)} \textbf{6} (1982), 9--24.
\href{https://doi.org/10.1090/S0273-0979-1982-14958-8}
{doi:10.1090/S0273-0979-1982-14958-8}.

\bibitem{Purcell}
J.~S.~Purcell, An introduction to fully augmented links,
\emph{Contemp.\ Math.} \textbf{541} (2011), 205--220.
\href{https://users.monash.edu.au/~jpurcell/papers/expos-auglinks.pdf}
{Author's manuscript}.

\bibitem{Rivin}
I.~Rivin, Euclidean structures on simplicial surfaces and hyperbolic
volume, \emph{Ann.\ of Math.} \textbf{139} (1994), 553--580.
\href{https://doi.org/10.2307/2118572}{doi:10.2307/2118572}.
\end{thebibliography}
